%% ======================================================================
\section{Preliminaries and Problem Formulation}
\label{sec:prelim}
%% ======================================================================

\paragraph{Sampling and candidates.}
Given a problem $x$ with reference answer $y$, a solver produces $n = 4$ chains
of thought and the answer each one reaches. Grouping answers that agree under a
type-aware normalisation yields $K \le n$ distinct \emph{candidates}
$c_1, \dots, c_K$ with vote counts $v_1, \dots, v_K$ summing to $n$, each
carrying one representative solution. Self-consistency \cite{wang2022self}
returns the candidate with the most votes, which we write $m(x)$ and refer to
throughout as CoT maj@4, our baseline. We write CoT pass@4 for the oracle that
is correct whenever $y$ appears among the candidates.

\paragraph{The selection task.}
We study \emph{answer selection}: given the problem and the candidates the
solver has already produced, commit to one of them. Every method we compare
consumes the same four samples, drawn once per instance and reused unchanged,
and differs only in the evidence it reads off them. CoT maj@4 uses the vote
counts; Universal Self-Consistency \cite{chen2023universal} sees the samples
with duplicates shown, so multiplicity remains visible; the verification
methods we study see the problem text and the representative solutions with
vote counts withheld. Holding the input fixed in this way makes the comparison
one about selection evidence rather than about sampling luck.

\subsection{Requirements on a gold-free selector}
\label{sec:requirements}

Three constraints separate a selector deployable at inference time from one
that merely explains a benchmark after the fact.

\begin{description}
\item[\textnormal{(R1) \emph{Gold-free.}}] The selector reads only the problem
  and the samples. The reference answer enters at scoring and nowhere else.
\item[\textnormal{(R2) \emph{No re-solving.}}] The selector issues no further
  call to the solver, so its cost is a fraction of the sampling budget it
  augments rather than a multiple of it.
\item[\textnormal{(R3) \emph{Non-destructive.}}] Because the selector may
  overturn $m(x)$, it must be judged by what it repairs \emph{and} what it
  damages, not by accuracy alone.
\end{description}

R3 fixes our unit of measurement. Against the maj@4 baseline, a \emph{fix} is
an instance the majority got wrong and the selector gets right, a \emph{break}
the reverse, and we report the two together with an exact McNemar test rather
than an accuracy delta on its own.

\paragraph{Selection is a two-sided bet.}
A selector constrained to the candidates cannot invent an answer the solver
never produced, which bounds both of its tails. Let $\mathcal{F}$ be the
instances where the majority is wrong but some candidate is right, and
$\mathcal{B}$ those where the majority is already right and the candidates
disagree. Over $N$ instances,
\begin{equation}
  -\frac{|\mathcal{B}|}{N}
  \;\le\; \Delta\mathrm{acc}
  \;\le\; \frac{|\mathcal{F}|}{N}
  \;=\; \mathrm{pass@}4 - \mathrm{maj@}4 ,
  \label{eq:bounds}
\end{equation}
since a unanimous-but-wrong instance contains no correct candidate and so
raises neither bound. The upper bound is the \emph{headroom} a selector
competes for. The two sets are not the same size, and we report both in
Section~\ref{sec:results}; the relevant point here is that R3 is forced by
Eq.~\eqref{eq:bounds} rather than adopted out of caution.

\paragraph{When a verdict is worth acting on.}
The vote is itself an informative signal, so a selector that overturns it needs
evidence strong enough to overcome that prior. Let $p(v)$ be the probability
that a candidate with $v$ votes is correct and
$\rho(v) = p(v)/(1 - p(v))$ the corresponding odds. A verifier returning
evidence $E$ about a candidate carries the likelihood ratio
$\Lambda = \Pr(E \mid \text{correct}) / \Pr(E \mid \text{wrong})$, and
preferring a $v_c$-vote candidate over a $v_m$-vote majority requires
\begin{equation}
  \Lambda \;>\; \rho(v_m) \,/\, \rho(v_c) .
  \label{eq:threshold}
\end{equation}
We estimate $p(v)$ on a development split, so Eq.~\eqref{eq:threshold} states
in advance how strong a checker must be before a verification call can pay for
itself. It is a specification fixed before any verifier is measured, not a
post-hoc account of one, and Section~\ref{sec:results} compares measured
likelihood ratios against it.

\paragraph{Perturbation pairs.}
Rather than evaluate on problems whose failure modes are unknown, we evaluate
on matched pairs. Each item is a frozen pair $(P, P')$ together with the
solution strategy $S$ that the original invites, where $P'$ alters one
structure, relationship or constraint of $P$ while remaining well posed and
keeping the original answer format. In a \emph{strategy-breaking} pair,
applying $S$ unchanged to $P'$ gives a wrong answer, which we record as the
\emph{trap}; in a \emph{strategy-preserving} pair, $S$ still works but the
answer changes and must be recomputed. The preserving mode is the control for
mere novelty, separating a method that degrades on any altered problem from one
that detects strategy misfit. The trap is what makes a pair informative: a
candidate equal to the trap is wrong for a known reason, which supplies
per-instance ground truth for the error class that frequency-based selection
cannot see. Section~\ref{sec:setup} gives the generation and verification
protocol and the resulting dataset statistics.
